\section{Introduction}
$\label{sec:Introduction}$
Memory management remains a primary bottleneck in modern high-performance computing (HPC) and distributed cluster environments. Large-scale scientific workflows across domains such as genomics, remote sensing, and materials science routinely submit thousands of concurrent tasks to resource managers like Slurm~\cite{yoo2003slurm}, HTCondor~\cite{thain2005htcondor}, or Kubernetes~\cite{burns2016borg}. Because cluster schedulers require declarative memory reservations prior to execution, accurate prelaunch resource estimation is critical. Underestimating memory requirements triggers out-of-memory (OOM) termination, causing cascading execution failures, costly checkpoint restarts, and broken dependencies across parallel pipeline stages~\cite{witt2019feedback,lehmann2024ponder}. Conversely, overestimating requirements leads to substantial resource overprovisioning, low cluster utilization, and increased queuing delays for concurrent workloads~\cite{bader2024sizey,witt2019feedback}. Because workflow developers cannot easily quantify the memory footprints of black-box analytical tools across highly variable datasets, reservations gravitate toward conservative upper bounds, resulting in cluster inefficiency~\cite{bader2024sizey,lehmann2024ponder}.

To automate resource scaling, machine learning models are increasingly deployed to predict task-level memory boundaries. Prior literature focuses predominantly on model selection, establishing that expressive non-linear architectures, such as $k$-nearest neighbors, random forests, and gradient-boosted trees—outperform traditional linear regressors at capturing complex runtime behavior~\cite{rodrigues2016helping,bader2024sizey}. However, predictive accuracy remains fundamentally bounded by the available feature space rather than model capacity. 


Current memory estimators suffer from two key information limitations.
First, standard schedulers rely predominantly on static prelaunch features, such as staged input file sizes, which fail to reflect actual runtime execution patterns. Existing predictors use aggregate file byte volumes, process names, and thread counts to estimate peak allocation thresholds~\cite{witt2019feedback,tovar2018jobsizing,bader2024sizey,lehmann2024ponder}. However, static storage metrics do not map to the exact data subsets parsed during execution. For instance, selective query tools read minor fractions of a large file index, whereas memory-mapped binaries load file segments into physical memory via page faults without issuing conventional \(\texttt{read}\) system calls. Consequently, resource allocators conflate overall staged storage volume with live resident memory requirements. Resolving this discrepancy requires capturing fine-grained runtime data consumption, which is computationally expensive to trace continuously and causally unavailable at the moment of scheduling.

 

%In genomics alignment pipelines~\cite{li2009bwa}, a task may open a massive reference genome for localized random access; aggregate storage statistics report gigabytes scanned, but peak resident set size (RSS) is driven entirely by the specific index subsets paged into physical memory. Without fine-grained, task-attributed consumption telemetry, predictors conflate overall staged storage volume with live memory footprints (Figure~\ref{fig:scatter}). Resolving this requires tracking kernel-level execution data, which is computationally expensive to collect and causally hidden at scheduling time.

Second, streaming workloads display high variance in historical data density, creating an asymmetry in feature specificity. Frequently executed workflow steps possess dense historical records of nearly identical past runs, whereas novel or updated pipeline stages lack empirical baselines. Monolithic global models degrade under this variance: they over-generalize well-characterized tasks by smoothing out tight, process-specific historical signals, yet generate over-optimistic or under-constrained bounds for unfamiliar tasks. An effective framework must isolate and route tasks dynamically according to the granularity of available empirical evidence.

This paper presents \toolname{}, a closed-loop memory-allocation framework designed to address these feature limitations. \toolname{} utilizes an eBPF-based auditing daemon to dynamically capture explicit read bytes and file-backed memory-mapped (\texttt{mmap}) page faults, directly attributing these I/O operations to individual task process trees. 
Completed tasks populate an asynchronous database, supplying peak RSS records globally and kernel-level data consumption metrics for instrumented runs. To exploit historical data at the appropriate granularity, \toolname{} structures empirical observations into hierarchical compatibility \textit{scopes}, ranging from specific instance repetitions to global hardware fallbacks. 

%Tasks are automatically routed to the strictest scope meeting minimum sample support requirements, while multi-scale summaries from adjacent scopes are preserved as features.
%CAMP maps these features through a multi-quantile gradient boosting architecture that predicts the upper tail of peak memory. Because true consumption is a runtime feature hidden at launch, CAMP deploys an intermediate proxy model to estimate prelaunch consumption metric. This estimate is appended to the feature space of a downstream quantile model, whose predictions are adjusted by an out-of-fold calibration layer using historical log-residuals. This calibration bounds over-prediction error while actively correcting for the under-coverage risks typical of data-sparse scopes.  %It implements an intermediate proxy model to estimate the unobservable prelaunch consumption memory value, combined with an out-of-fold calibration layer that adjusts predictions based on historical residual distribution limits. The resulting allocation request reflects a calibrated safety margin rather than unweighted statistical averages. 

To bound tracing overhead on production clusters, \toolname{} operates under a strict administrative profiling budget, deploying a multi-gate active learning selection policy. These gates prioritize tasks with a high probability of allocation failure, high anticipated training information utility, or those operating within under-observed feature regions. The resulting system evaluates a continuous loop of multi-quantile prediction, selective auditing, kernel measurement, and asynchronous retraining.
% This work addresses the following research questions:
% \begin{itemize} 
% \item \textbf{RQ1:} How do strace and eBPF mechanisms differ regarding data-consumption coverage and online execution tracing overhead? 
% \item \textbf{RQ2:} Under a unified model family, data partition, and evaluation protocol, how does the inclusion of historical and predicted runtime consumption metrics affect allocation safety and cluster capacity reservation costs? 
% \item \textbf{RQ3:} For an equivalent profiling budget, do fixed active-learning selection gates identify higher-risk or more informative tasks than uniform or stratified random sampling? 
% \item \textbf{RQ4:} Can the integrated cold-start, persistence, retraining, and fresh-execution feedback loop operate end-to-end within a live production environment? 
% \end{itemize}
The primary contributions of this work are: 
\begin{itemize}
%\item the design, implementation, and empirical evaluation of \toolname{}, an end-to-end, closed-loop %resource management 
%framework that introduces task-attributed data consumption telemetry to parallel workflow scheduling; 
\item a \textit{scope}-based hierarchical prediction architecture combining multi-scale history summarization and non-parametric quantile ensembles; %global consumption proxy estimators, and residual tail calibration layers; 
\item  an evaluation of fixed multi-gate selective auditing policies against uniform and process-stratified random baselines under equivalent sampling budgets; 
 \item  a comprehensive experimental evaluation on several Bioinformatics workflows including comparative performance baseline against Sizey~\cite{bader2024sizey} and a live end-to-end Bowtie2 case study evaluating the framework across cold-start, persistence, model updating, and execution phases.
\end{itemize}

To the best of our knowledge this is the first time a distribution based model is used for peak memory prediction. Past work~\cite{bader2024sizey} has predominantly considered point-based models. 

%We evaluate these systems components comprehensively, analyzing structural trade-offs alongside performance gains: comprehensive eBPF tracing introduces non-negligible instrumentation overhead in the current prototype, maximizing execution safety guarantees shifts reservation profiles, and point-prediction error metrics do not change uniformly across all workloads.

The remainder of this paper is structured as follows. Section~\ref{sec:problem} formalizes the problem. Section~\ref{sec:design} details the architectural design of \toolname{}, including the hierarchical scope definitions, history construction, multi-quantile inference, and active selection policies. Section~\ref{sec:experiments} defines the experimental workloads, evaluation protocol, baselines, and tracking metrics. Section~\ref{sec:results} discusses the empirical results, Section~\ref{sec:related} covers related literature, and Section~\ref{sec:conclusion} concludes the paper.
